% ---
% filename : "electrotechnique_machines_20260923_418_nb_paires_poles_glissement.tex"
% id : "electrotechnique_machines_20260923_418"
% title : "Nombre de paires de pôles et glissement d'un moteur asynchrone"
% domain : "Électrotechnique"
% subdomain : "Machines"
% tags : ["moteur_asynchrone", "triphasé", "glissement", "frequence_synchronisme", "paires_de_pôles"]
% difficulty : 1
% time_solve : 8
% has_octave : true
% octave_files : ["electrotechnique_machines_20260923_418_nb_paires_poles_glissement.m"]
% author : "om"
% ---
\documentclass[a4paper,12pt,answers,addpoints]{exam}

% ---------------------------------------------------------
% LANGUE, ENCODAGE, FONTS
% ---------------------------------------------------------
\usepackage[utf8]{inputenc}

% ---------------------------------------------------------
% GÉOMÉTRIE ET MISE EN PAGE
% ---------------------------------------------------------
\usepackage[left=1.7cm,right=1.5cm,top=2cm,bottom=1cm]{geometry}
\usepackage{microtype}
\usepackage{float}
\usepackage{needspace}

% ---------------------------------------------------------
% MATHÉMATIQUES
% ---------------------------------------------------------
\usepackage{amsmath, amssymb, bm, esvect}

\newcommand{\V}{\overrightarrow}
\def\vect#1{\overrightarrow{\strut#1}}
\providecommand{\Degrees}[0]{\ensuremath{^\circ}}

% ---------------------------------------------------------
% UNITÉS ET GRANDEURS (SIunitx)
% ---------------------------------------------------------
\usepackage{siunitx}
\sisetup{output-decimal-marker={,}}
\DeclareSIUnit \var {var}
\DeclareSIUnit \VA {VA}
\DeclareSIUnit \kVA {kVA}
\DeclareSIUnit[quantity-product={}] \degree{\text{\textdegree}}

% Permet la césure dans les nombres avec unités (évite les Underfull \hbox)
%\AllowBreakAtQQ

% ---------------------------------------------------------
% GRAPHIQUES : TikZ + CircuitikZ (sans tkz-fct qui cause des erreurs spath3)
% ---------------------------------------------------------
\usepackage{tikz}
\usetikzlibrary{
	arrows.meta,
	intersections,
	decorations.pathreplacing,
	%	calligraphy,
	positioning
}

\usepackage[european,straightvoltages,EFvoltages]{circuitikz}
\usepackage{pgfplots}
\pgfplotsset{compat=1.12}

% Symbole étoile-triangle personnalisé
\newcommand{\wye}{%
	\mathbin{\tikz[x=1ex,y=1ex]{%
			\draw[line width=.1ex] (0,0)--(45:1)--++(-45:1) (45:1)--++(0,1);
	}}%
}

% ---------------------------------------------------------
% TABLEAUX
% ---------------------------------------------------------
\usepackage{tabularx}
\usepackage{tabu}
\usepackage{longtable}

% ---------------------------------------------------------
% HYPERLIENS
% ---------------------------------------------------------
\usepackage[colorlinks=true,
menucolor=red,
breaklinks=true,
urlcolor=blue,
linkcolor=blue]{hyperref}

% ---------------------------------------------------------
% COMMANDES PERSONNELLES
% ---------------------------------------------------------
\newcommand\nomauteur{bdminasmorgul@protonmail.com}
\newcommand\dateTe{23.09.2026}
\newcommand\brancheTe{TS}

% ---------------------------------------------------------
% STYLE DE L'EXERCICE
% ---------------------------------------------------------
\pagestyle{headandfoot}
\firstpageheadrule
\runningheadrule

% En-tête avec largeur fixe pour éviter les dépassements
\firstpageheader{\brancheTe\ (\nomauteur)}{Élève : }{\makebox[3.5cm][r]{\dateTe\ \thepage/\numpages}}
\runningheader{\brancheTe\ (\nomauteur)}{Élève : }{\makebox[3.5cm][r]{\dateTe\ \thepage/\numpages}}

% Pied de page
\newcommand{\totalpointtts}{%
	\begin{tikzpicture}[remember picture, overlay]
		\node[anchor=south west, inner sep=0pt, outer sep=0pt]
		at ([xshift=0.5cm,yshift=0.5cm]current page.south west)
		{\rotatebox{90}{Total des points : \numpoints}};
	\end{tikzpicture}%
}

\firstpagefooter{\hspace*{\fill}\totalpointtts}{}{}
\runningfooter{\hspace*{\fill}\totalpointtts}{}{}

% Réglages exam
\printanswers
\addpoints
\pointsinmargin
\bracketedpoints

% ---------------------------------------------------------
% LISTINGS (code)
% ---------------------------------------------------------
\usepackage{xcolor}
\usepackage{listings}

\lstdefinestyle{octavestyle}{
	language=Matlab,
	basicstyle=\ttfamily\small,
	keywordstyle=\color{blue},
	commentstyle=\color{green!50!black},
	stringstyle=\color{red},
	stepnumber=1,
	frame=single,
	breaklines=true,
	breakatwhitespace=true,
	tabsize=4
}

\usepackage{enumitem}
\setlist[enumerate,1]{label=\alph*), ref=\alph*)}
\newenvironment{explication}{%
	\par\noindent\textbf{Explication. }\itshape
}{\par\normalfont}
\renewcommand\partlabel{\arabic{partno}.}
\begin{document}
	\unframedsolutions
	\pointsinmargin
	\bracketedpoints
	
\section*{Préparation TS PQ CFC ELMO,IELE,PELE}
\begin{questions}
\question
Un moteur asynchrone triphasé est relié au réseau \qty{400}{[V]}/\qty{230}{[V]} – \qty{50}{[Hz]}. En pleine charge, sa fréquence de rotation est \qty{1425}{[tr/min]}.

Déterminer :

\begin{parts}
    \part[3] le nombre de paires de pôles du stator.
    \part[3] le glissement du moteur en pleine charge.
\end{parts}

\begin{solution}

La fréquence de synchronisme d'un moteur asynchrone triphasé dépend directement du nombre de paires de pôles selon la relation suivante, où $f$ est la fréquence du réseau et $p$ le nombre de paires de pôles~:
\begin{align*}
    n_s = \dfrac{60 \cdot f}{p}
\end{align*}
Pour $f = \qty{50}{[Hz]}$, les fréquences de synchronisme possibles selon le nombre de paires de pôles sont~:
\begin{align*}
    p = 1 \quad &\Rightarrow \quad n_s = \qty{3000}{[tr/min]} \\[6pt]
    p = 2 \quad &\Rightarrow \quad n_s = \qty{1500}{[tr/min]} \\[6pt]
    p = 3 \quad &\Rightarrow \quad n_s = \qty{1000}{[tr/min]} \\[6pt]
    p = 4 \quad &\Rightarrow \quad n_s = \qty{750}{[tr/min]}
\end{align*}

1. Nombre de paires de pôles

On exprime $p$ à partir de la relation de synchronisme :
\begin{align*}
    p = \dfrac{60 \cdot f}{n_s} = \dfrac{60 \cdot 50}{1500} = \dfrac{3000}{1500} = 2
\end{align*}
Le stator possède donc \boxed{2 \text{ paires de pôles}}

2. Glissement en pleine charge

Le glissement est défini par le rapport entre la fréquence de glissement ($n_s - n$) et la fréquence de synchronisme~:
\begin{align*}
    s = \dfrac{n_s - n}{n_s} = \dfrac{1500 - 1425}{1500} = \dfrac{75}{1500} = 0,05 \\[6pt]
    s_{\%} = 0,05 \cdot 100 = 5 \text{ \%}
\end{align*}

Le glissement en pleine charge est donc \boxed{s = 0,05 = 5 \text{ \%}}, valeur typique pour un moteur asynchrone à rotor à cage en régime nominal.
\newpage
\begin{verbatim}
% Télécharger OCTAVE depuis https://wiki.octave.org/Using_Octave et l'installer
% Puis copier-coller le code dans OCTAVE

% Données de l'exercice
f = 50;        % Fréquence du réseau [Hz]
n = 1425;      % Fréquence du rotor en pleine charge [tr/min]

% 1. Recherche du nombre de paires de pôles
% Les fréquences de synchronisme possibles sont n_s = 60*f/p
% On recherche p tel que n_s est la valeur la plus proche de n
% (le rotor tourne toujours légèrement en dessous de n_s)
p = 1;
while (true)
  n_s = 60 * f / p;
  printf("p = %d paires de pôles  ->  n_s = %.0f tr/min\n", p, n_s);
  if (n_s <= n)
    % On conserve la dernière valeur de n_s supérieure à n
    if (p > 1)
      p = p - 1;
      n_s = 60 * f / p;
    end
    break;
  end
  p = p + 1;
end

printf("\n--- Résultats ---\n");
printf("Nombre de paires de pôles du stator : p = %d (soit %d pôles)\n", p, 2*p);
printf("Fréquence de synchronisme retenue     : n_s = %.0f tr/min\n", n_s);

% 2. Calcul du glissement
g = (n_s - n) / n_s;
printf("Glissement                          : g = %.4f\n", g);
printf("Glissement en pourcentage           : g = %.1f %%\n", g * 100);
\end{verbatim}

\end{solution}
\end{questions}
\end{document}